\documentclass[12pt,a4paper]{ctexart}

\usepackage[top=2.5cm, bottom=2.5cm, left=2.5cm, right=2.5cm]{geometry}
\usepackage{amsmath,amssymb}
\usepackage{booktabs}
\usepackage{enumitem}
\usepackage[colorlinks=true,linkcolor=blue,urlcolor=blue]{hyperref}
\usepackage{setspace}
\onehalfspacing

\title{模型预测结果文档说明（Introduction）}
\author{TeamBDATA2600409}
\date{2026年8月30日}

\begin{document}
\maketitle

\section{设计思路}

\subsection{问题概述}

本赛题基于 50,000 条电动汽车消费者记录，完成三个回归任务，评分指标为调整决定系数（adj-$R^2$）：

\begin{table}[h]
\centering
\begin{tabular}{ccc}
\toprule
\textbf{任务} & \textbf{预测目标} & \textbf{最终模型} \\
\midrule
T1 月度能耗 & monthly\_energy\_consumption\_kwh & $0.8163487 \times weekly$ \\
T2 月度充电成本 & monthly\_charging\_cost & $0.8163487 \times weekly \times price$ \\
T3 燃油替代支出 & fuel\_expense\_per\_month & $8.3947643 \times daily$ \\
\bottomrule
\end{tabular}
\end{table}

\subsection{核心思路：物理结构优先，复杂模型克制}

本方案的可行性建立在数据探索阶段的一个关键发现之上：该数据集由近似确定性的物理关系生成，稀疏公式即可逼近预测上限。具体证据如下：

\begin{enumerate}[itemsep=0pt]
\item \textbf{比率诊断}：月能耗与周里程的比率（$energy/weekly$）近似一个与其余特征无关的随机量，各车型差异不足 0.001；燃油支出与日通勤过原点 OLS 得系数 8.395，残差标准差约 40；充电成本与``能耗 $\times$ 电价''的 $R^2$ 达 0.9994。
\item \textbf{相关性结构}：月能耗与周里程相关系数 0.936，燃油支出与日通勤 0.952，三者目标与对应出行变量呈极强的线性关系。
\item \textbf{复杂模型证伪}：在稀疏基线的折内残差上训练浅层 LightGBM，三个任务相对基线的 $\Delta R^2$（Bootstrap 95\% 置信区间）全部为负，题面建议的车型/地域交互特征亦无增益。
\end{enumerate}

据此，本方案以可解释的闭式公式为最终模型，其优势在于：与数据生成机制一致、无过拟合风险、在 adj-$R^2$ 评分口径下不受自由度惩罚、且完全可复现。

\subsection{关键设计：严格嵌套 OOF 级联（T2）}

任务二的充电成本近似等于``能耗 $\times$ 电价''。为避免在训练时使用真实能耗、测试时改用预测能耗造成的 teacher forcing（训练/推理分布错位），本方案采用严格嵌套 OOF 级联：在任务二的外层训练折内部再次交叉验证任务一，为训练行生成内层 OOF 能耗，仅以该预测值构造 $cost = \hat{energy} \times price$，从机制上杜绝跨折信息泄漏。

\subsection{异常值处理（T3）}

燃油支出存在 271 个负值。消融实验表明，保留原值（过原点 OLS）优于删除、缩尾与 Huber 稳健回归，故最终模型保留原始标签、不做对数变换与预测截断。

\section{环境\&语言}

\begin{table}[h]
\centering
\begin{tabular}{ll}
\toprule
\textbf{项目} & \textbf{配置} \\
\midrule
语言 & Python 3.11 \\
包管理器 & uv \\
依赖库 & pandas, numpy, scipy, scikit-learn, lightgbm \\
可选依赖 & catboost, xgboost, optuna, shap, matplotlib \\
算力 & CPU（无需 GPU），全部实验数十秒内完成 \\
操作系统 & Windows 11 \\
随机种子 & 42 / 2026 / 7 \\
数据规模 & 50,000 行 $\times$ 23 列 \\
\bottomrule
\end{tabular}
\end{table}

\section{模型调用方式}

\subsection{最终模型（闭式公式）}

最终模型不依赖任何序列化文件，参数仅 2 个系数（存档于 \texttt{artifacts/models/final\_models.json}）：

\begin{equation}
\hat{y}_1 = 0.8163487 \times weekly\_travel\_distance\_km
\end{equation}
\begin{equation}
\hat{y}_2 = 0.8163487 \times weekly\_travel\_distance\_km \times electricity\_cost\_per\_kwh
\end{equation}
\begin{equation}
\hat{y}_3 = 8.3947643 \times daily\_commute\_km
\end{equation}

\subsection{复现与预测}

解压支撑材料后，在项目根目录执行：

\begin{verbatim}
uv sync                              # 安装依赖
uv run python -m src.schema_check    # 数据校验
uv run python -m src.folds           # 生成固定折
uv run python -m src.experiments     # 基线 + 消融
uv run python -m src.finalize        # 残差确认 + 最终确认
uv run python -m src.submit          # 拟合最终模型 + 生成提交
\end{verbatim}

对正式测试集预测：

\begin{verbatim}
uv run python -m src.submit --test <测试.csv>
\end{verbatim}

输出三份提交文件 \texttt{T1/T2/T3\_submission.csv}，格式为 \texttt{row\_id, true, prediction}。数据 CSV 需置于 \texttt{复赛B题/B题数据集.csv}。

\subsection{指标（5 折 $\times$ 3 种子 pooled OOF）}

\begin{table}[h]
\centering
\begin{tabular}{cccc}
\toprule
\textbf{任务} & \textbf{adj-$R^2$} & \textbf{RMSE} & \textbf{MAE} \\
\midrule
T1 & 0.876298 & 31.04 & 24.59 \\
T2 & 0.917817 & 7.15 & 5.32 \\
T3 & 0.905416 & 40.03 & 31.96 \\
\bottomrule
\end{tabular}
\end{table}

\end{document}
